\section{Models and implementation}
\label{sec:methods}
\subsection{B0 and B0-UW: centralized controls}
B0 learns full $64$-dimensional user and item embeddings centrally using BPR,
with Gaussian initialization of SD $0.01$. Its final retained configuration is
Adam learning rate $10^{-3}$, explicit L2 coefficient $10^{-2}$, batch size
1,024, at most 1,000 epochs, and patience of 100 validation checks. B0-UW uses
the same modeling and optimization settings but samples a user uniformly and
then a positive uniformly within that user. B0 samples interactions uniformly.
The matched three-seed experiment separates the change in user weighting from
the change to federated optimization more clearly than B0--B1 alone.

Earlier B0 documents retain superseded learning rates and stopping rules.
The report uses the later common-stopping-protocol results as the final
reference. Historical search tables remain part of the audit trail rather
than being blended into a new average.

\subsection{B1: federated full matrix without DP}
In each B1 round, users are independently included with $q=0.1$. A selected
user starts from the current shared $Q$ and its persistent local vector,
then performs two local epochs on its own positives. The local SGD learning
rate is $5$, regularization is $10^{-5}$, and the server learning rate is $1$.
The large numerical difference from B0's Adam rate reflects a different
optimizer and the normalization of each local loss; it is not a directly
comparable step-size scale. Selected users' item differences are averaged
with equal client weight. B1 divides by realized participation and skips
empty rounds because it has no private-release requirement.

Validation occurs every ten rounds. Patience is 100 validation checks, with
an 8,000-round maximum, and the best validation checkpoint is restored for
evaluation. B1 is the converged nonprivate federated reference. Local user
vectors are not transmitted, but neither federation nor local storage supplies
an additional formal guarantee in this nonprivate baseline.

\subsection{B2: private full matrix}
B2 applies \cref{eq:aggregate} to the complete shared item update while retaining
B1's local training settings. Its final private horizon is 50 rounds with
$C=1$ and server rate $1$. It always evaluates the final round, without
privacy-dependent early stopping. The fixed denominator, full-coordinate
noise, and empty-round accounting distinguish it from a simple ``add noise
to B1'' implementation.

The earlier B2 protocol used 1,000 rounds, $C=1.5$, and server rate $2$.
Validation studies motivated the final 50-round protocol. Three settings
changed together, so the improvement cannot be attributed uniquely to
shorter training. The retained no-DP control at the 50-round horizon is
necessary: the gap from a converged B1 includes a training-budget effect
even when no noise is present.

\subsection{B3 and B4: jointly trained item factors}
B3 represents $Q=AB$ at ranks $4,8,16,32$ and trains both factors without DP.
User vectors retain dimension $64$. Regularization is applied to $p_u$ and
the effective item rows $A_iB$, with gradients through both factors. Adding
independent L2 penalties on $A$ and $B$ would change the objective and the
role of scale, so that is not done.

QR decompositions and a core SVD rebalance the factors while preserving their
product up to floating-point tolerance. This avoids an extreme reciprocal
scaling of $A$ and $B$, but it does not remove a perturbation already present
in $AB$. It also introduces a changing orthogonal coordinate convention,
which becomes important in the later trajectory probes. The final B3 local
rates are $0.625,1.25,1.25,0.625$ in increasing rank order. B3 follows B1's
convergence rule with one disclosed exception: rank 32, seed 123 uses a
predeclared 16,000-round diagnostic ceiling and stops at 9,350, with best
checkpoint 8,350. All other final runs stop below the original ceiling.

B4 applies a single joint clip to $(\Delta A,\Delta B)$, then independent
Gaussian coordinates across that concatenated vector. It uses the same
$q,T,C,\sigma,\delta$ as final B2. Its final rank-4 local/server rates are
$2.5/1$; the other ranks use $5/0.5$. Historical B4 received more validation
tuning than B2, including a bounded fallback after unstable candidate
trajectories. The B2--B4 comparison is therefore a comparison of retained
methods, not a clean intervention on coordinate count alone.

\subsection{E1: a separately controlled comparison}
E1 preserves B0--B4 and creates three explicit families: Full, Two, and FixedB.
Full repeats the final full-matrix mechanism; Two trains both factors; FixedB
holds a public row-orthonormal $B$ fixed and privately trains only $A$.
The basis is seeded independently of interactions (public seed 314159).
At each rank, Two and FixedB begin with the same effective $Q$. This removes
an avoidable initialization discrepancy between the two factorized methods.

Each of nine units (Full and two methods at four ranks) receives four
local/server learning-rate candidates: $(1.25,1)$, $(2.5,1)$, $(5,0.5)$,
and $(5,1)$. Selection uses three seeds at reference privacy target four,
subject to finite behavior across six levels: no clipping/no noise,
clipping only, and the four private targets. The selected local rate is
$5$ throughout. The server rate is $1$ for Full and FixedB, $0.5$ for
Two ranks 4/8/16, and $1$ for Two rank 32. FixedB's local rate reaches the
upper grid boundary, a limitation on claims about its optimal attainable utility.

The fixed public basis is useful analytically. Let $v_u=Bp_u$. Then
$s_{ui}=A_i^Tv_u$, and every $v_u\in\R^r$ can be realized by
$p_u=B^Tv_u$ when $BB^T=I_r$. Thus FixedB retains the same hard score-rank
bound as Two at rank $r$; its failure cannot be attributed simply to a
smaller maximum score rank. Jointly learning $B$ changes the latent metric,
gradient dynamics, clipping, and aggregation behavior. It does not increase
the maximum rank beyond $r$. \Cref{app:theory} gives the equivalence precisely.

\subsection{User-level private popularity}
For each user define a binary indicator $h_u\in\{0,1\}^M$ for distinct
training-positive items, and bound its contribution by
\begin{equation}
 x_u=h_u\min\left(1,\frac{\sqrt{20}}{\norm{h_u}_2}\right),\qquad
 z=\sum_u x_u+\mathcal N(0,20\sigma_{\mathrm{pop}}^2I).
 \label{eq:pop}
\end{equation}
The sum has add/remove user sensitivity $\sqrt{20}$. A single full-population
release is calibrated analytically at the same target $(\epsilon,\delta)$
as the iterative models, using its own multiplier. Items are ranked by $z_i$,
with the same user-specific held-out candidate filter. This is a useful
nonpersonalized collaborative reference because it spends its budget on one
bounded statistic instead of repeated iterative updates.

The no-noise bounded-popularity score is not the same as raw item counts:
high-activity users are downweighted by the norm bound. The historical raw
popularity, bounded no-noise popularity, and private popularity must therefore
be named separately. A local metadata-only model and a public-features-plus-
private-residual model were proposed as future directions but were not trained.

\subsection{Communication accounting}
The simulator counts dense float32 uploads and downloads for participating
clients, excluding local $P$. At a client-round, the modeled payload is
\begin{equation}
 b_{\mathrm{Full}}=8Md,\qquad
 b_{\mathrm{Two}}=8r(M+d),\qquad
 b_{\mathrm{FixedB}}=8Mr\quad\text{bytes}.
 \label{eq:comm}
\end{equation}
The factor eight is four bytes per scalar in each direction. FixedB can
reconstruct its public basis from a seed; explicit provision would add a
one-time $4rd$ bytes per client. Network framing, compression, cryptographic
overhead, failures, and real device latency are not measured. Total modeled
traffic accumulates the payload over realized selected clients and rounds.
It is distinct from payload size and from utility under a fixed total budget.
